\documentclass[12pt,]{book}
\usepackage{lmodern}
\usepackage{setspace}
\setstretch{1}
\usepackage{amssymb,amsmath}
\usepackage{ifxetex,ifluatex}
\usepackage{fixltx2e} % provides \textsubscript
\ifnum 0\ifxetex 1\fi\ifluatex 1\fi=0 % if pdftex
  \usepackage[T1]{fontenc}
  \usepackage[utf8]{inputenc}
\else % if luatex or xelatex
  \ifxetex
    \usepackage{mathspec}
  \else
    \usepackage{fontspec}
  \fi
  \defaultfontfeatures{Ligatures=TeX,Scale=MatchLowercase}
\fi
% use upquote if available, for straight quotes in verbatim environments
\IfFileExists{upquote.sty}{\usepackage{upquote}}{}
% use microtype if available
\IfFileExists{microtype.sty}{%
\usepackage{microtype}
\UseMicrotypeSet[protrusion]{basicmath} % disable protrusion for tt fonts
}{}
\usepackage[marginpar=2cm, top=3cm, bottom=4cm]{geometry}
\usepackage{hyperref}
\hypersetup{unicode=true,
            pdfborder={0 0 0},
            breaklinks=true}
\urlstyle{same}  % don't use monospace font for urls
\usepackage{natbib}
\bibliographystyle{apalike}
\usepackage{longtable,booktabs}
\usepackage{graphicx,grffile}
\makeatletter
\def\maxwidth{\ifdim\Gin@nat@width>\linewidth\linewidth\else\Gin@nat@width\fi}
\def\maxheight{\ifdim\Gin@nat@height>\textheight\textheight\else\Gin@nat@height\fi}
\makeatother
% Scale images if necessary, so that they will not overflow the page
% margins by default, and it is still possible to overwrite the defaults
% using explicit options in \includegraphics[width, height, ...]{}
\setkeys{Gin}{width=\maxwidth,height=\maxheight,keepaspectratio}
\IfFileExists{parskip.sty}{%
\usepackage{parskip}
}{% else
\setlength{\parindent}{0pt}
\setlength{\parskip}{6pt plus 2pt minus 1pt}
}
\setlength{\emergencystretch}{3em}  % prevent overfull lines
\providecommand{\tightlist}{%
  \setlength{\itemsep}{0pt}\setlength{\parskip}{0pt}}
\setcounter{secnumdepth}{5}
% Redefines (sub)paragraphs to behave more like sections
\ifx\paragraph\undefined\else
\let\oldparagraph\paragraph
\renewcommand{\paragraph}[1]{\oldparagraph{#1}\mbox{}}
\fi
\ifx\subparagraph\undefined\else
\let\oldsubparagraph\subparagraph
\renewcommand{\subparagraph}[1]{\oldsubparagraph{#1}\mbox{}}
\fi

%%% Use protect on footnotes to avoid problems with footnotes in titles
\let\rmarkdownfootnote\footnote%
\def\footnote{\protect\rmarkdownfootnote}

%%% Change title format to be more compact
\usepackage{titling}

% Create subtitle command for use in maketitle
\newcommand{\subtitle}[1]{
  \posttitle{
    \begin{center}\large#1\end{center}
    }
}

\setlength{\droptitle}{-2em}
  \title{}
  \pretitle{\vspace{\droptitle}}
  \posttitle{}
  \author{}
  \preauthor{}\postauthor{}
  \date{}
  \predate{}\postdate{}

% Améliore l'esthétique de la police
\usepackage{lmodern}

%Packages pour créer des tableaux 
\usepackage{longtable} % Pour des tableaux dont la longueur dépasse une feuille A4
\usepackage{tabularx} % Pour des tableaux à largeur définie
\usepackage{array} % Pour améliorer la qualité typographique des tableaux.

\usepackage{amsthm}
\newtheorem{theorem}{Theorem}[chapter]
\newtheorem{lemma}{Lemma}[chapter]
\theoremstyle{definition}
\newtheorem{definition}{Definition}[chapter]
\newtheorem{corollary}{Corollary}[chapter]
\newtheorem{proposition}{Proposition}[chapter]
\theoremstyle{definition}
\newtheorem{example}{Example}[chapter]
\theoremstyle{definition}
\newtheorem{exercise}{Exercise}[chapter]
\theoremstyle{remark}
\newtheorem*{remark}{Remark}
\newtheorem*{solution}{Solution}
\begin{document}

%Page de garde
\begin{titlepage}
\frontmatter
\begin{figure}[t]
\includegraphics[width=5cm]{../figures-ext/LogoParisDescartes}
\end{figure}

\begin{center}
UNIVERSITÉ PARIS DESCARTES \\
\vspace*{1cm}
\textbf{École doctorale}\\
\vspace*{0,5cm}
\textit{Laboratoire/équipe de recherche}\\
\vspace*{1cm}
\LARGE{\textbf{Titre de la thèse}}\\
\vspace*{0,5cm}
\large{\textit{\textbf{Sous-titre de la thèse}}}\\
\vspace*{2cm}
\large{\textbf{Par [Prénom et nom de l'auteur]}}\\
\vspace*{1cm}
Thèse de doctorat de [Discipline : consulter la liste des disciplines]\\
\vspace*{1cm}
Dirigée par [Prénom et nom du directeur de thèse]\\
\vspace*{1cm}
\small{Présentée et soutenue publiquement le [date de soutenance]}\\
\end{center}
\vspace*{1cm}
\begin{footnotesize}
Devant un jury composé de : \\
\begin{tabular}{lll}
Prénom NOM & [fonction] - université\\
Prénom NOM & [fonction] - université\\
Prénom NOM & [fonction] - université\\
\end{tabular}
\end{footnotesize}

\begin{figure}[b]
\begin{center}
\includegraphics{../figures-ext/creativecommons}
\end{center}
\end{figure}




\clearpage


%Abstract
\newpage
\thispagestyle{empty}
\noindent % Supprime le retrait de paragraphe
\textbf{Résumé (français) :}
\vskip 1cm
\noindent
\textbf{Title :}
\vskip 1cm
\noindent
\textbf{Abstract :}
\vskip 1cm
\noindent
\textbf{Mots-clés (français) :}
\vskip 1cm
\noindent
\textbf{Keywords :}


%Dédicace
\newpage
\emph{Dédicace}
\vspace*{\fill}
 \begin{quote}
 \emph{Suffering has been stronger than all other teaching, and has taught me to understand what your heart used to be. I have been bent and broken, but - I hope - into a better shape.}
 \end{quote}
 \vspace*{\fill}

%Remerciements
\newpage
\thispagestyle{empty}
\begin{center}
\large{\textbf{Avertissement}}
\end{center}
\vspace{2cm}
Cette thèse de doctorat est le fruit d’un travail approuvé par le jury de soutenance et
réalisé dans le but d’obtenir le diplôme d’Etat de docteur de philosophie. Ce document
est mis à disposition de l’ensemble de la communauté universitaire élargie.
Il est soumis à la propriété intellectuelle de l’auteur. Ceci implique une obligation de
citation et de référencement lors de l’utilisation de ce document.
D’autre part, toute contrefaçon, plagiat, reproduction illicite encourt toute poursuite
pénale.
\vspace*{\fill}

\emph{Code de la Propriété Intellectuelle. Articles L 122.4 \newline
Code de la Propriété Intellectuelle. Articles L 335.2-L 335.10}


\newpage
\thispagestyle{empty}
\begin{center}
\large{\textbf{Remerciments}}
\end{center}
\vspace{2cm}


\end{titlepage}

{
\setcounter{tocdepth}{4}
\tableofcontents
}
\listoftables
\listoffigures
\hypertarget{intro}{%
\chapter{Immuno-biology of cancer}\label{intro}}

\setcounter{page}{11}\renewcommand{\thepage}{\arabic{page}}\vspace*{\fill}

\begin{quote}
 \emph{Suffering has been stronger than all other teaching, and has taught me to understand what your heart used to be. I have been bent and broken, but - I hope - into a better shape.}
 \end{quote}

\vspace*{\fill}

\begin{quote}
\emph{Suffering has been stronger than all other teaching, and has
taught me to understand what your heart used to be. I have been bent and
broken, but - I hope - into a better shape.}

--- someone
\end{quote}

This chapter will introduce basic topic of cancer and participation of
stroma in cancer development, progression and response to treatment. It
will also describe

\hypertarget{cancer-seen-as-complex-environment}{%
\section{Cancer seen as complex
environment}\label{cancer-seen-as-complex-environment}}

For a long time studying tumor was focused on tumor cells, their
reporogramming, mutations. It was seen as diseas of uncontrolled cells.
Recent reseach moved research focus from tumor cells to tumor cells in
their proper context : tumor micorenvironment.

\hypertarget{our-understanding-of-cancer-over-time}{%
\subsection{Our understanding of cancer over
time}\label{our-understanding-of-cancer-over-time}}

cancer is a disease touching blah blah many ppl over the word. it has
been known that blah blah and then types

\hypertarget{tumor-micro-environment-fiend-or-foe}{%
\subsection{Tumor micro environment : fiend or foe
?}\label{tumor-micro-environment-fiend-or-foe}}

what is tme : composition, roles it was decided the environment is bad
for cancer Tumors effectively supresses immune response : activates
negative regulatory pathways ( checkpoints)

For ages we didn't know much about how modulate tme Now we know it can
do both - review hallmarks of cancer immuno

\hypertarget{cancer-immune-phenotypes}{%
\subsection{Cancer immune phenotypes}\label{cancer-immune-phenotypes}}

There can be distinguished cancer phenotypes depending on immune
infiltration how they are measure, defined, indexes, types of cancer,
impact

\begin{quote}
In further support of a role for memory T cells in antitumour responses,
tumour-infiltrating lymphocytes that express CD4 or CD8 extracted from
experimental tumour models typically have the features of memory T cells
and can possess an activated or exhausted phenotype, expressing markers
such as PD-1, T-cell immunoglobulin and mucin-domain containing protein
3 (TIM-3) and lymphocyte activation gene 3 (LAG-3). (IMMUNE CANCER
CIRCLE)
\end{quote}

\begin{quote}
Anticancer immunity in humans can be segregated into three main
phenotypes: the immune-desert phenotype (brown), the immune--excluded
phenotype (blue) and the inflamed phenotype (red). (IMMUNE CANCER CIRCLE
Fig 3)
\end{quote}

\hypertarget{immune-signatures}{%
\subsection{Immune signatures}\label{immune-signatures}}

definition of signature: marker genes, list of genes, weighted list we
can talk about general immune signature of signature of immune
infiltration and stroma or immune signature of a specific cell type of
functional subpopulation purpose of signatures

avaliability of immune signatures

problem of non cosistence of immune signatures origin of signatures

\hypertarget{immunotherapies}{%
\section{Immunotherapies}\label{immunotherapies}}

This section outlines progress in cancer therapies with a focus on
immune therapies. It will link the ongoing research on TME with
therapeutical potential.

\hypertarget{cancer-therapies}{%
\subsection{Cancer therapies}\label{cancer-therapies}}

\hypertarget{recent-progress-in-immuno-therapies}{%
\subsection{Recent progress in
immuno-therapies}\label{recent-progress-in-immuno-therapies}}

most potential\\
cytotoxic T-lymphocyte protein 4 (CTLA4) and programmed cell death
protein 1 (PD-1)

\begin{quote}
CTLA4 is a negative regulator of T cells that acts to control T-cell
activation by competing with the co-stimulatory molecule CD28 for
binding to shared ligands CD80 (also known as B7.1) and CD86 (also known
as B7.2).The cell-surface receptor PD-1 is expressed by T cells on
activation during priming or expansion and binds to one of two ligands,
PD-L1 and PD-L2. Many types of cells can express PD-L1, including tumour
cells and immune cells after exposure to cytokines such as interferon
(IFN)-γ; however, PD-L2 is expressed mainly on dendritic cells in normal
tissues. Binding of PD-L1 or PD-L2 to PD-1 generates an inhibitory
signal that attenuates the activity of T cells. The `exhaustion' of
effector T cells was identified through studies of chronic viral
infection in mice in which the PD-L1/PD-1 axis was found to be an
important negative feedback loop that ensures immune homeostasis; it is
also an important axis for restricting tumour immunity. (IMMUNE CANCER
CIRCLE)
\end{quote}

\begin{quote}
The mechanisms that underlie cancer immunotherapy differ considerably
from those of other approaches to cancer treatment. Unlike chemotherapy
or oncogene-targeted therapies, cancer immunotherapy relies on promoting
an anticancer response that is dynamic and not limited to targeting a
single oncogenic derangement or other autonomous feature of cancer
cells. Cancer immunotherapy can therefore lead to antitumour activity
that simultaneously targets many of the abnormalities that differentiate
cancer cells and tumours from normal cells and tissues.(IMMUNE CANCER
CIRCLE)
\end{quote}

\hypertarget{potential-of-developpement-of-new-immunotherapies}{%
\subsection{Potential of developpement of new
immunotherapies}\label{potential-of-developpement-of-new-immunotherapies}}

\begin{quote}
As effective as immunotherapy can be, only a minority of people exhibit
dramatic responses, with the frequency of rapid tumour shrinkage from
single-agent anti-PD-L1/PD-1 antibodies ranging from 10--40\%, depending
on the individual's indication ( Zou, W., Wolchok, J. D. \& Chen, L.
PD-L1 (B7-H1) and PD-1 pathway blockade for cancer therapy: mechanisms,
response biomarkers, and combinations. Sci. Transl. Med. 8, 328rv4
(2016).)
\end{quote}

\begin{quote}
predicting reponse: The immune-inflamed phenotype correlates generally
with higher response rates to anti-PD-L1/PD-1 therapy51,62,67,69,70,71,
which suggests that biomarkers could be used as predictive tools. Most
attention has been paid to PD-L1, which is thought to reflect the
activity of effector T cells because it can be adaptively expressed by
most cell types following exposure to IFN-γ6,82. IMMUNE CANCER CIRCLE)
\end{quote}

\hypertarget{methods}{%
\chapter{Mathematical foundation of deconvolution}\label{methods}}

Here is a review of existing methods.

\hypertarget{deconvolution-of-transcriptomes-and-methylomes}{%
\chapter{Deconvolution of transcriptomes and
methylomes}\label{deconvolution-of-transcriptomes-and-methylomes}}

We describe our methods in this chapter.

\hypertarget{comparative-analysis-of-cancer-immune-infiltration}{%
\chapter{Comparative analysis of cancer immune
infiltration}\label{comparative-analysis-of-cancer-immune-infiltration}}

Some \emph{significant} applications are demonstrated in this chapter.

\hypertarget{example-one}{%
\section{Example one}\label{example-one}}

\hypertarget{example-two}{%
\section{Example two}\label{example-two}}

\hypertarget{heterogeneity-of-immune-cell-types}{%
\chapter{Heterogeneity of immune cell
types}\label{heterogeneity-of-immune-cell-types}}

We have finished a nice book.

\hypertarget{annexes-annexes}{%
\chapter*{(ANNEXES) Annexes}\label{annexes-annexes}}
\addcontentsline{toc}{chapter}{(ANNEXES) Annexes}

\bibliography{book.bib,packages.bib}


\end{document}
